\documentclass[11pt]{article}
\usepackage{varnernotes}

\newcommand{\R}{\mathbb{R}}

\begin{document}

\lecturenotes{9b}{Option Contracts: Call and Put Payoff and Profit}{Fall 2026}{October 22, 2026}{Jeffrey D. Varner}

\learningobjectives
  {Read an option contract}
  {Identify the underlying, strike, expiration, exercise style, premium, position direction, and contract multiplier.}
  {Separate payoff from profit}
  {Derive long and short call and put expiration payoffs, premium-adjusted profit, breakeven prices, and loss bounds.}
  {Compose option positions}
  {Represent a multi-leg strategy as a signed sum of contract cash flows without hiding units or financing conventions.}

\section{A Derivative Is a Contract on Another Quantity}

A derivative is a financial contract whose cash flows depend on an underlying quantity such as an equity price, interest rate, currency, commodity, index, or credit event. Forwards and futures obligate counterparties to transact later; swaps exchange specified cash-flow streams; options give the holder a right while imposing a corresponding obligation on the writer if the right is exercised.

An option description is incomplete without its terms. Let $S_t\geq0$ be the underlying price in USD/share at time $t$, let $K>0$ be the strike in USD/share, and let $T>0$ be expiration. A \emph{call} gives its long holder the right to buy the underlying at $K$; a \emph{put} gives its long holder the right to sell at $K$. European-style exercise occurs only at $T$, while American-style exercise may occur on or before $T$. This lecture analyzes expiration cash flows; early-exercise decisions require a time-dependent valuation model.

The premium is paid when the position opens. Let $C_0\geq0$ and $P_0\geq0$ be call and put premiums in USD/share at time $0$. Let $q\in\mathbb N$ be the deliverable multiplier in shares/contract. Standard U.S. equity options commonly use $q=100$, but adjusted and nonstandard contracts can differ; the actual contract specification controls.

\begin{figure}[ht]
  \centering
  \begin{tikzpicture}[>=Latex,font=\sffamily\small,node distance=0.7cm]
    \node[draw=vnink,fill=vnpanel,minimum width=2.1cm,minimum height=0.8cm] (terms) {contract terms};
    \node[draw=vnrule,fill=vnsoft,right=of terms,minimum width=2.1cm,minimum height=0.8cm] (state) {realized $S_T$};
    \node[draw=vnrule,fill=vnsoft,right=of state,minimum width=2.1cm,minimum height=0.8cm] (payoff) {expiration payoff};
    \node[draw=vnrule,fill=vnsoft,right=of payoff,minimum width=2.1cm,minimum height=0.8cm] (profit) {premium-adjusted P\&L};
    \draw[vnlink,thick,->] (terms)--(state);
    \draw[vnlink,thick,->] (state)--(payoff);
    \draw[vnlink,thick,->] (payoff)--(profit);
    \node[below=0.18cm of terms,text=vnmuted] {$K,T$, style, $q$};
    \node[below=0.18cm of profit,text=vnmuted] {plus financing/cost convention};
  \end{tikzpicture}
  \caption{Payoff is fixed by the contract and the terminal underlying price. Profit additionally depends on the opening premium, position direction, multiplier, financing convention, fees, and any intervening exercise or assignment.}
  \label{fig:contract-to-profit}
\end{figure}

\section{Expiration Payoffs Are Contractual}

For any real $x$, write $x^+:=\max(x,0)$. The per-share terminal payoff to a long European call and long European put is
\begin{equation}
  \label{eq:long-payoffs}
  c_T(S_T):=(S_T-K)^+,
  \qquad
  p_T(S_T):=(K-S_T)^+.
\end{equation}
The call is exercised at expiration when buying at $K$ is better than buying at $S_T$; the put is exercised when selling at $K$ is better than selling at $S_T$. The holder may decline exercise, so both payoffs are nonnegative.

Let $\theta\in\{-1,+1\}$ denote position direction, with $+1$ long and $-1$ short. Ignoring counterparty default, the writer's terminal payoff is the negative of the holder's payoff:
\begin{equation}
  \label{eq:signed-payoff}
  V_T^{c,\theta}=\theta c_T(S_T),
  \qquad
  V_T^{p,\theta}=\theta p_T(S_T).
\end{equation}
This payoff symmetry does not imply symmetric risk after the premium: call and put terminal functions have different bounds because $S_T\geq0$ but has no finite contractual upper bound.

\begin{definition}{Moneyness at expiration}{moneyness}
A call is in the money when $S_T>K$, at the money when $S_T=K$, and out of the money when $S_T<K$. A put reverses the inequalities. Moneyness describes exercise value relative to the strike; it does not by itself determine whether the original position earned a profit because the premium was paid or received separately.
\end{definition}

\section{Profit Requires a Cash-Flow Convention}

The common classroom expiration diagram subtracts the opening premium without growing it through time. Under this undiscounted cash-flow convention, per-share profit is
\begin{equation}
  \label{eq:profit-convention}
  \Pi_T^{c,\theta}(S_T)=\theta\bigl(c_T(S_T)-C_0\bigr),
  \qquad
  \Pi_T^{p,\theta}(S_T)=\theta\bigl(p_T(S_T)-P_0\bigr).
\end{equation}
For a long position, the premium is an initial outflow; for a short position, it is an initial inflow. Equation~\eqref{eq:profit-convention} is useful for shape and breakeven calculations, but it mixes time-$0$ and time-$T$ dollars. A time-$T$ economic P\&L would subtract the premium's financed future value, while a time-$0$ net present value would discount the terminal payoff. Fees and bid--ask effects must also be added consistently.

Under the classroom convention, the long call breakeven solves $S_T-K-C_0=0$ in the exercised region, and the long put breakeven solves $K-S_T-P_0=0$. Therefore
\begin{equation}
  \label{eq:breakevens}
  S_{\mathrm{BE}}^{c}=K+C_0,
  \qquad
  S_{\mathrm{BE}}^{p}=K-P_0.
\end{equation}
The same prices are breakevens for the corresponding short positions because their profits are exact negatives under \eqnref{profit-convention}. If $P_0>K$, the algebraic put breakeven is negative and therefore outside the feasible domain $S_T\geq0$.

\begin{figure}[ht]
  \centering
  \begin{tikzpicture}[x=0.78cm,y=0.72cm,>=Latex,font=\sffamily\small]
    \begin{scope}
      \draw[vnink,->] (0,0)--(6.4,0) node[right] {$S_T$};
      \draw[vnink,->] (0,-1.5)--(0,3.1) node[above] {USD/share};
      \draw[vnlink,line width=1.5pt] (0,0)--(3.2,0)--(6.0,2.8);
      \draw[vncarnelian,dashed,line width=1.4pt] (0,-0.65)--(3.2,-0.65)--(6.0,2.15);
      \draw[vnmuted,dotted] (3.2,-1.3)--(3.2,2.8) node[above] {$K$};
      \node[vnink] at (2.7,2.55) {long call};
      \node[text=vnlink] at (5.25,2.75) {payoff};
      \node[text=vncarnelian] at (5.35,1.55) {profit};
    \end{scope}
    \begin{scope}[xshift=8.4cm]
      \draw[vnink,->] (0,0)--(6.4,0) node[right] {$S_T$};
      \draw[vnink,->] (0,-1.5)--(0,3.1) node[above] {USD/share};
      \draw[vnlink,line width=1.5pt] (0,2.8)--(3.2,0)--(6.0,0);
      \draw[vncarnelian,dashed,line width=1.4pt] (0,2.15)--(3.2,-0.65)--(6.0,-0.65);
      \draw[vnmuted,dotted] (3.2,-1.3)--(3.2,2.8) node[above] {$K$};
      \node[vnink] at (2.7,2.55) {long put};
      \node[text=vnlink] at (0.75,2.75) {payoff};
      \node[text=vncarnelian] at (0.8,1.55) {profit};
    \end{scope}
  \end{tikzpicture}
  \caption{Long call and put payoffs (solid blue) and classroom profit diagrams (dashed red). The premium shifts profit downward without changing the kink at $K$; breakeven is not the same point as the strike.}
  \label{fig:payoff-profit}
\end{figure}

\section{Loss Bounds and Direction Matter}

For one share-equivalent under the convention in \eqnref{profit-convention}, the terminal bounds are summarized below. Multiply every amount by the deliverable $q$ and by the number of contracts to obtain position-level dollars.

\begin{center}
\begin{tabular}{@{}llll@{}}
  \textbf{Position} & \textbf{Maximum profit} & \textbf{Maximum loss} & \textbf{Breakeven} \\
  \hline
  Long call  & unbounded as $S_T\to\infty$ & $C_0$ & $K+C_0$ \\
  Short call & $C_0$ & unbounded as $S_T\to\infty$ & $K+C_0$ \\
  Long put   & $K-P_0$ at $S_T=0$ & $P_0$ & $K-P_0$ \\
  Short put  & $P_0$ & $K-P_0$ at $S_T=0$ & $K-P_0$ \\
\end{tabular}
\end{center}

The short put's terminal loss is large but finite when the underlying price is bounded below by zero. Calling it ``unlimited downside loss'' is incorrect for an ordinary put with finite strike. By contrast, an uncovered short call has no finite contractual loss bound because the underlying price has no finite upper bound. Margin requirements, early assignment, liquidity, and forced liquidation can make realized risk severe well before terminal bounds are reached.

\notebooklink
  {L9b: Single-Contract Payoff and Profit}
  {https://github.com/varnerlab/CHEME-5660-CourseRepository-Fall-2026/blob/main/lectures/week-9/L9b/CHEME-5660-L9b-Example-SingleContractPayoffProfit-Fall-2025.ipynb}
  {Build long and short call and put contracts from an option chain, compute expiration payoff and premium-adjusted profit, and verify the analytical breakevens.}

\section{Composite Positions Add Leg by Leg}

Let a strategy contain legs $j\in\{1,\ldots,d\}$. For leg $j$, let $\theta_j\in\{-1,+1\}$ be direction, $n_j\in\mathbb N$ the number of contracts, $q_j\in\mathbb N$ the deliverable multiplier, $V_{j,T}(S_T)$ the long per-unit payoff, and $\mathcal P_{j,0}\geq0$ the per-unit premium. Under the same undiscounted convention, total terminal profit is
\begin{equation}
  \label{eq:composite-profit}
  \Pi_{\mathrm{strategy},T}(S_T)
  =\sum_{j=1}^{d}\theta_jn_jq_j
  \bigl[V_{j,T}(S_T)-\mathcal P_{j,0}\bigr].
\end{equation}
Linearity makes vertical spreads, straddles, strangles, covered calls, and protective puts transparent: their kinks occur at the component strikes, while their vertical offsets depend on the net opening debit or credit. Exercise style, expiration, underlying, deliverable, and settlement must be compatible before legs are casually combined into one diagram.

\notebooklink
  {L9b: Composite-Contract Payoff and Profit}
  {https://github.com/varnerlab/CHEME-5660-CourseRepository-Fall-2026/blob/main/lectures/week-9/L9b/CHEME-5660-L9b-Example-CompositeContractsPL-Fall-2025.ipynb}
  {Combine signed option legs to study vertical spreads, straddles, and strangles, and identify their piecewise profit regions and breakevens.}

\keytakeaways
  {In this lecture, we translated option contract terms into expiration cash flows and kept contractual payoff distinct from the investor's premium-dependent profit.}
  {Payoff belongs to the contract}
  {A long call pays $(S_T-K)^+$ and a long put pays $(K-S_T)^+$; short positions reverse those cash flows.}
  {Profit needs a convention}
  {Premium, financing time, multiplier, fees, and position direction determine P\&L, so strike, moneyness, and breakeven are different concepts.}
  {Loss bounds are asymmetric}
  {An uncovered short call has unbounded terminal loss, while a standard short put's maximum terminal loss is finite at $S_T=0$.}
  {Next, value European claims with Black--Scholes--Merton and American claims with a Cox--Ross--Rubinstein exercise lattice.}

\begin{readinglist}
  \item John C. Hull, \emph{Options, Futures, and Other Derivatives}, 11th ed. (2022), chapters on option-market mechanics and payoff strategies.
  \item Options Clearing Corporation, \href{https://www.theocc.com/company-information/documents-and-archives/options-disclosure-document}{\emph{Characteristics and Risks of Standardized Options}}, for contract, exercise, assignment, and risk disclosures.
  \item Robert C. Merton, \href{https://doi.org/10.2307/1831029}{``Theory of Rational Option Pricing,''} \emph{Bell Journal of Economics and Management Science} 4(1), 141--183 (1973).
\end{readinglist}

\vfill
{\sffamily\footnotesize\color{vnmuted}\textbf{Educational use and risk notice.} These notes are for instruction only and do not constitute investment advice. Options involve substantial risk, and short positions can create obligations exceeding the premium received.}

\end{document}
